\subsection{表示向测度空间的延拓}

本节中假设 K = C。

设 \(\varrho\) 是 G 在分离局部凸拟完备空间 E 中的连续线性表示。记 \(\mathcal{M}_{c}(\mathrm{G})\) 为 G 上具有紧支撑且赋予紧收敛拓扑的测度空间；它是空间 \(\mathcal{C}(\mathrm{G})\) 的对偶。对于每个测度 \(\nu \in \mathcal{M}_{c}(\mathrm{G})\)，置

\[
\varrho (\nu) = \int_ {\mathrm{G}} \varrho (g) d \nu (g).
\]

这是 \(\mathcal{L}(\mathrm{E})\) 的一个元素。特别地，\(\varrho(\varepsilon_g) = \varrho(g)\) 对所有 \(g \in \mathrm{G}\) 成立。

根据 INT, VIII, p. 136, § 2, no. 6，映射 \(\nu \mapsto \varrho(\nu)\) 是从 \(\mathcal{M}_c(\mathrm{G})\) 到赋予紧收敛拓扑的空间 \(\mathcal{L}(\mathrm{E})\) 的连续线性映射。根据 INT, VIII, p. 145, § 3, no. 3, prop. 11，它是含单位代数的态射。

设 \(x \in \mathrm{E}\)。当 \(\mathcal{M}_c(\mathrm{G})\) 赋予紧收敛拓扑时，由 \(\nu \mapsto \varrho(\nu)x\) 定义的从 \(\mathcal{M}_c(\mathrm{G})\) 到 E 的映射是连续的（INT, VI, p. 27, § 1, no. 7, prop. 16 以及 EVT, III, p. 31, prop. 4）。

对于在 G 的某个紧子集之外为零的 \(f \in \mathcal{L}^{1}(\mathrm{G})\)，测度 \(f \cdot \mu\) 具有紧支撑；我们把 E 的自同态 \(\varrho^{\mu}(f)\)（也记为 \(\varrho(f)\)）记为 \(\varrho(f \cdot \mu)\)。此自同态只取决于其类 \(\widetilde{f}\)，即 \(f\) 在 \(\mathrm{L}^{1}(\mathrm{G})\) 中的类，我们也把它记为 \(\varrho^{\mu}(\widetilde{f})\) 或 \(\varrho(\widetilde{f})\)。

同样，设 \(\varrho\) 是 G 在巴拿赫空间 E 中的连续有界线性表示。对于每个有界测度 \(\nu \in \mathcal{M}^{1}(\mathrm{G})\)，置

\[
\varrho (\nu) = \int_ {\mathrm{G}} \varrho (g) d \nu (g).
\]

根据 INT, VIII, 同上，映射 \(\nu \mapsto \varrho(\nu)\) 是从 G 上有界复测度代数 \(\mathcal{M}^{1}(G)\) 到巴拿赫代数 \(\mathcal{L}(E)\) 的连续含单位巴拿赫代数态射。

令 \(\rho\) 为 G 上的函数 \(g\mapsto\|\varrho(g)\|\)；INT, VIII, p. 145 的命题 11 中记作 \(\mathcal{M}^{\rho}\) 的代数（其元素是测度 \(\nu\)，满足 \(\rho\in\mathcal{L}^{1}(\nu)\)）与巴拿赫代数 \(\mathcal{M}^{1}(G)\) 重合。事实上，置 \(M = \sup \rho\)。有 \(M > 0\)，因为 \(\rho(e)=1\)，且 \(M^{-1} \leqslant \rho \leqslant M\)，因为不等式 \(\|\varrho(e)\|\leqslant\|\varrho(g)\|\|\varrho(g^{-1})\|\) 对所有 \(g\in G\) 成立。因此，\(\rho\in\mathcal{L}^{1}(\nu)\) 当且仅当 \(\nu\) 是有界测度。

若 \(f \in \mathcal{L}^{1}(\mathrm{G})\)，自同态 \(\varrho^{\mu}(f)\)（或简记为 \(\varrho(f)\)）就是 \(\varrho(f \cdot \mu)\)；对于类 \(\tilde{f}\)，它是 \(f\) 在 \(\mathrm{L}^{1}(\mathrm{G})\) 中的类，同样如此。

引理 1. --- 设 \(\varrho\) 是 G 在希尔伯特空间 E 中的酉表示。映射 \(\nu \mapsto \varrho(\nu)\) 从 \(\mathcal{M}^{1}(\mathrm{G})\) 到 \(\mathcal{L}(\mathrm{E})\)，是对合巴拿赫代数的含单位态射。

由前述内容，只需证明态射 \(\nu \mapsto \varrho(\nu)\) 是对合的。设 \(\nu \in \mathcal{M}^{1}(\mathrm{G})\)。测度 \(\nu^{*}\) 是测度 \(\check{\nu}\) 的共轭测度（I, p. 99, 例 4）。根据共轭测度的定义（INT, III, p. 52, § 1, no. 5），计算得

\[
\begin{array}{c} \langle x \mid \varrho (\nu) y \rangle = \int_ {\mathrm{G}} \langle x \mid \varrho (g) y \rangle d \nu (g) \\ = \int_ {\mathrm{G}} \langle \varrho (g ^ {- 1}) x \mid y \rangle d \nu (g) = \int_ {\mathrm{G}} \langle \varrho (g) x \mid y \rangle d \check {\nu} (g) \\ = \overline {{\int_ {\mathrm{G}} \langle y \mid \varrho (g) x \rangle d \nu^ {*} (g)}} = \langle \varrho (\nu^ {*}) x \mid y \rangle \end{array}
\]

对 E 中所有 \(x\) 和 \(y\) 成立，从而 \(\varrho(\nu)^{*} = \varrho(\nu^{*})\)。

设 \(\varrho\) 是 G 在局部凸拟完备空间 E（分别在巴拿赫空间 E）中的连续线性表示（分别为连续有界表示）。如果 F 是 E 的闭子空间，并且子表示 \(\varrho_{\mathrm{F}}\) 由 \(\varrho\) 在 F 上诱导，那么对于每个测度 \(\nu \in \mathcal{M}_c(\mathrm{G})\)（分别为 \(\nu \in \mathcal{M}^1 (\mathrm{G})\)），线性映射 \(\varrho (\nu)\) 保持子空间 F 不变，并且在 F 上与 \(\varrho_{\mathrm{F}}(\nu)\) 重合。

反之，设 \(F \subset E\) 是闭子空间，并且在每个线性映射 \(\varrho(f)\) 下稳定，其中 \(f \in \mathcal{H}(G)\)。则 F 定义 \(\varrho\) 的一个子表示（INT, VIII, p. 139, § 2, no. 7, prop. 10）。

回忆（A, VIII, p. 388），函数 \(f \colon \mathbf{G} \to \mathbf{C}\) 称为中心函数，是指对所有 \(g\) 和 \(h\)（均属于 \(\mathbf{G}\)），有 \(f(gh) = f(hg)\)。这等价于说 \(f\) 在共轭作用下不变。

定义 1. --- 有界测度 \(\nu \in \mathcal{M}^{1}(\mathrm{G})\) 称为中心的，是指 \(\varepsilon_{g} * \nu = \nu * \varepsilon_{g}\) 对所有 \(g \in \mathrm{G}\) 成立。

如果 G 是单模群，且 \(f \in \mathcal{C}(\mathrm{G})\) 关于 \(\mu\) 可积，那么测度 \(f \cdot \mu\) 中心，当且仅当 f 是中心函数。

设 \(\varrho\) 是 G 在巴拿赫空间 E 中的连续有界表示（分别是 G 在拟完备分离局部凸空间 E 中的连续表示）。对于 G 上的每个有界中心测度 \(\nu\)（分别为 G 上的任意紧支撑中心测度 \(\nu\)），线性映射 \(\varrho(\nu)\) 是从 \(\varrho\) 到 \(\varrho\) 的 G-态射。事实上，对每个 \(g \in G\)，有

\[
\varrho (\nu) \varrho (g) = \varrho (\nu * \varepsilon_ {g}) = \varrho (\varepsilon_ {g} * \nu) = \varrho (g) \varrho (\nu).
\]

